% sec1b.tex — c03｜原文 p6–7｜§1.2 Global symmetry
% 编号公式 (1.18)–(1.30) 共 13 条；编号环境：习题 1.1。
% [原文 p.6]
这里，
\[
\begin{aligned}
i_B(x) &= \Tr_{(\text{bosonic letters with } H=0)} x^H ,\\
i_F(x) &= \Tr_{(\text{fermionic letters with } H=0)} x^H .
\end{aligned}
\tag{1.18}
\]
\begin{exercise}[1.1]
证明这一点。
\end{exercise}
对超对称谐振子，玻色单字母指标 $i_B(x) = 0$，费米单字母指标 $i_F(x) = 0$。因此 $i(x) = 0$。这给出完整的指标为
\[
I(x) = PE[i(x)] = 1 .
\tag{1.19}
\]
重复一遍：对相互作用的超对称理论，配分函数不易计算，特别地它是 $x$ 的函数；而指标容易计算（把耦合设为零即可），并且不依赖于 $x$（因为只有 $H=0$ 的态有贡献）。

\subsection{整体对称性}

考虑上述拉氏量的 $2N$ 个副本
\[
L = \frac{1}{2}|\dot x_i|^2 - \frac{1}{2}|x_i|^2 + \ii\bar\psi_i\dot\psi_i - \bar\psi_i\psi_i .
\tag{1.20}
\]
它具有 $\SO(2N)$ 对称性。有一种引入保持该对称性的相互作用的方法：
\[
L = \frac{1}{2}\dot x_i^2 - \frac{1}{2}\left|\frac{\partial W(x_i)}{\partial x_i}\right|^2 + \ii\bar\psi_i\dot\psi_i - \frac{\partial^2 W(x)}{\partial x_i\partial x_j}\bar\psi_i\psi_j .
\tag{1.21}
\]
取 $W = |x_i|^2$ 便得到自由理论。眼下让我们专注于自由理论。由于理论是自由的，配分函数很容易计算。它包含 $2N$ 个玻色产生算符 $a_{B,i}^\dagger$ 与 $2N$ 个费米产生算符 $a_{F,i}^\dagger$。这意味着
\[
z_B(x) = 2N x ,\qquad z_F(x) = 2N x .
\tag{1.22}
\]
利用齐次指数公式，
\[
Z(x) = \left(\frac{1+x}{1-x}\right)^{2N} .
\tag{1.23}
\]
有时记录态在整体对称性下的表示是有用的。在此情形下，我们具有 $\SO(2N)$ 对称性。作法是引入相对于 Cartan 生成元的逸度。
\[
Z(x,a_i) = \Tr x^H a_i^{J_i} ,\qquad i = 1,\ldots,N .
\tag{1.24}
\]
% [原文 p.7]
最好用单字母配分函数来计算它。由于玻色子处在 $\SO(N)$ 的基础表示中。它们的单字母配分函数为
\[
z_B(x,a_i) = x\left(a_1 + \frac{1}{a_1} + \ldots + a_n + \frac{1}{a_n}\right) = x\,\chi_{\fund}(a_i) .
\tag{1.25}
\]
这里 $\chi_{\fund}(a_i)$ 是基础表示的特征标。一般地，一个表示的特征标定义为
\[
\chi_R(a_i) = \sum_{\rho\in R}\prod_i a_i^{\rho_i} \equiv \sum_{\rho\in R} a^\rho ,
\tag{1.26}
\]
其中 $\rho = (\rho_1,\ldots,\rho_N)$ 是权向量，求和遍历该表示的全部权向量。有时我们会把 $\prod_i a_i^{\rho_i} = a^\rho$ 缩写。回忆一下，伴随表示的非零权向量称为根。于是我们有，
\[
\begin{aligned}
z_B(x,a_i) &= x\left(a_1 + \frac{1}{a_1} + \ldots + a_n + \frac{1}{a_n}\right) = x\,\chi_{\fund}(a_i)\\
Z_B(x,a_i) &= PE[z(x,a)] = \prod_{i=1}^{N}\frac{1}{(1-x a_i)(1-x/a_i)}\\
z_F(x,a_i) &= x\left(a_1 + \frac{1}{a_1} + \ldots + a_n + \frac{1}{a_n}\right) = x\,\chi_{\fund}(a_i)\\
Z_F(x,a_i) &= \widetilde{PE}[z(x,a)] = \prod_{i=1}^{N}(1+x a_i)(1+x/a_i)\\
Z(x,a_i) &= Z_B(x,a_i)Z_F(x,a_i) = \prod_{i=1}^{N}\frac{(1+x a_i)(1+x/a_i)}{(1-x a_i)(1-x/a_i)} \equiv \prod_{\rho}\frac{(1+x a^\rho)}{(1-x a^\rho)} .
\end{aligned}
\tag{1.27}
\]
显然，当取 $a_i = 1$ 时，$Z(x,a_i)$ 退化为未精细化的配分函数 $Z(x)$。记录相对于 Cartan 生成元的荷的好处在于，我们可以计算具有给定整体对称性表示的态上的配分函数。这要用到特征标的正交性，
\[
\frac{1}{|W|}\oint\prod_{i=1}^{N}\frac{\dd a_i}{2\pi\ii a_i}\,\Delta(a_i)\chi_R(a_i)\chi_{R'}(a_i) = \delta_{RR'} .
\tag{1.28}
\]
其中 $N$ 是群的秩，$|W|$ 是相伴 Weyl 群的阶，$\Delta(a_i)$ 称为 Vandermonde 行列式。Vandermonde 行列式有下述用齐次指数表示的普遍公式\footnote{需要正则化才能使齐次指数良定义。}
\[
\Delta(a_i) = \left(PE\Big[\sum_{\alpha} a^\alpha\Big]\right)^{-1} = \prod_{\alpha}(1-a^\alpha) ,
\tag{1.29}
\]
% [原文 p.8]
其中 $\alpha = (\alpha_1,\ldots,\alpha_N)$ 是李代数的根。现在很清楚如何投影到具有给定表示的态上，
\[
Z(x,a_i)\big|_{R} = \frac{1}{|W|}\oint\prod_{i=1}^{N}\frac{\dd a_i}{2\pi\ii a_i}\,\Delta(a_i)Z(x,a_i)\chi_R(a_i) .
\tag{1.30}
\]
在计算规范理论的指标时，我们将需要投影到具有平庸表示的态上。平庸表示的特征标为 $\chi_R(a_i) = 1$。
% [Tutor Note: 术语：$i_B$、$i_F$ 是「单字母指标」，与 $z_B$、$z_F$（单字母配分函数）只差 $(-1)^F$ 分级；习题 1.1 要证的正是 (1.17) 在本例中给出 $I=1$。]
% [Tutor Note: 动机：$2N$ 个副本即 $N$ 个复谐振子，$\SO(2N)$ 来自把这 $2N$ 个实分量整体旋转。]
% [Tutor Note: 易错：(1.20) 用复变量 $x_i$ 写成 $|x_i|^2$，所以「$2N$ 个副本」指 $2N$ 个实自由度，不是 $2N$ 个复变量。]
% [Tutor Note: 易错：$\chi_R$ 是群论「特征标」，不要与 $\sI$（指标 index）混淆；本译本统一译作「特征标」。]
% [Tutor Note: 术语：$a^\rho$ 只是 $\prod_i a_i^{\rho_i}$ 的缩写，$\rho$ 是权向量而非单个整数。]
% [Tutor Note: 动机：加入 Cartan 逸度后 $Z(x,a_i)$ 的展开按 $\SO(2N)$ 表示分解；令全部 $a_i=1$ 即退回未精细化配分函数。]
% [Tutor Note: 易错：原文此处写 $\SO(N)$（前文为 $\SO(2N)$），而 (1.25) 的展开恰是 $\SO(2N)$ 矢量表示的特征标；疑为原文笔误，已登记 ERRATA，译文照录。]
% [Tutor Note: 文献：特征标正交性 (1.28) 是 Weyl 积分公式的投影形式；$\frac{1}{|W|}$ 与 Haar 测度配套出现。]
% [Tutor Note: 术语：$\Delta(a_i)=\prod_\alpha(1-a^\alpha)$ 遍历全部根（含正负根），故等于 Weyl 分母的绝对值平方；脚注 2 说明 PE 求逆需先正则化。]
% [Tutor Note: 易错：(1.29) 对 $\alpha$ 的乘积遍历全部根，不能只取正根，否则会漏掉 $1-a^{-\alpha}$ 因子。]
